% Chapter 9: Heterogeneity Analysis - Conditional Average Treatment Effects
% ============================================================================
% This chapter covers:
% - CATE estimation theory and practice
% - CausalForestDML for non-parametric heterogeneity
% - Best linear predictor of CATE
% - Heterogeneity visualization and interpretation
% - Policy-relevant subgroup analysis
% ============================================================================

\chapter{Heterogeneity Analysis}
\label{ch:heterogeneity}

\section{Introduction}\label{sec:ch09_introduction}

Chapters 1--8 focused on estimating the \textbf{Average Treatment Effect} (ATE)---a single number summarizing the causal effect across an entire population. But in practice, treatment effects often vary across individuals or units. Understanding this \textbf{heterogeneity} is essential for:

\begin{itemize}
    \item \textbf{Targeting}: Who benefits most from treatment?
    \item \textbf{Policy design}: Should treatment be universal or selective?
    \item \textbf{Mechanism understanding}: Why do effects vary?
    \item \textbf{External validity}: Will effects generalize to new populations?
\end{itemize}

\keyresult{
\textbf{Chapter Objectives:}
\begin{itemize}
    \item Define and estimate Conditional Average Treatment Effects (CATE)
    \item Apply CausalForestDML for non-parametric heterogeneity discovery
    \item Construct the Best Linear Predictor (BLP) for effect heterogeneity
    \item Visualize and interpret heterogeneous effects
    \item Connect to policy-relevant decision-making
\end{itemize}
}

\subsection{From ATE to CATE}

The Average Treatment Effect answers ``\textit{what is the effect on average?}'' The Conditional Average Treatment Effect answers ``\textit{what is the effect for units with characteristics $X = x$?}''

\begin{definition}[Conditional Average Treatment Effect]
\label{def:cate}
For a unit with covariates $X = x$, the CATE is:
\begin{equation}
\tau(x) = \mathbb{E}[Y(1) - Y(0) \mid X = x]
\label{eq:cate_definition}
\end{equation}
The ATE is the expectation of CATE over the covariate distribution:
\begin{equation}
\tau_{\text{ATE}} = \mathbb{E}[\tau(X)] = \mathbb{E}[\mathbb{E}[Y(1) - Y(0) \mid X]]
\end{equation}
\end{definition}

\insight{
\textbf{Why Heterogeneity Matters}

Consider a training program evaluation:
\begin{itemize}
    \item ATE = \$1,000 wage increase (statistically significant)
    \item But: $\tau(x) = \$3,000$ for workers without college degree
    \item And: $\tau(x) = -\$500$ for workers with graduate degrees
\end{itemize}

A universal program might be inefficient; targeted rollout could double the benefit at half the cost.
}

\subsection{Chapter Roadmap}

\begin{itemize}
    \item \textbf{Section \ref{sec:ch09_cate_theory}}: CATE Theory---identifiability and estimation challenges
    \item \textbf{Section \ref{sec:ch09_dml_heterogeneity}}: DML for Heterogeneous Effects---extending the framework
    \item \textbf{Section \ref{sec:ch09_causal_forest}}: CausalForestDML---non-parametric heterogeneity
    \item \textbf{Section \ref{sec:ch09_blp}}: Best Linear Predictor---summarizing heterogeneity
    \item \textbf{Section \ref{sec:ch09_visualization}}: Visualization and Interpretation
    \item \textbf{Section \ref{sec:ch09_subgroups}}: Subgroup Analysis and Policy Trees
    \item \textbf{Section \ref{sec:ch09_exercises}}: Exercises
\end{itemize}


% ============================================================================
\section{CATE Theory}\label{sec:ch09_cate_theory}
% ============================================================================

\subsection{Identification}

CATE identification requires the same assumptions as ATE, applied conditionally:

\begin{theorem}[CATE Identification]
\label{thm:cate_identification}
Under the following assumptions:
\begin{enumerate}
    \item \textbf{Conditional unconfoundedness}: $(Y(0), Y(1)) \perp T \mid X$
    \item \textbf{Overlap}: $0 < P(T = 1 \mid X = x) < 1$ for all $x$ in the support
    \item \textbf{SUTVA}: No interference, single version of treatment
\end{enumerate}
the CATE is identified as:
\begin{equation}
\tau(x) = \mathbb{E}[Y \mid T = 1, X = x] - \mathbb{E}[Y \mid T = 0, X = x]
\end{equation}
\end{theorem}

\subsection{Estimation Challenges}

CATE estimation faces challenges beyond ATE:

\begin{enumerate}
    \item \textbf{High dimensionality}: $\tau(x)$ is a function, not a scalar
    \item \textbf{Curse of dimensionality}: Local estimation becomes sparse
    \item \textbf{Regularization bias}: Standard ML methods induce bias in conditional means
    \item \textbf{Inference complexity}: Constructing valid CIs for $\tau(x)$ requires special care
\end{enumerate}

\begin{definition}[Nuisance Functions for CATE]
\label{def:cate_nuisance}
CATE estimation requires estimating:
\begin{align}
\mu_0(x) &= \mathbb{E}[Y \mid T = 0, X = x] \quad \text{(control outcome)} \\
\mu_1(x) &= \mathbb{E}[Y \mid T = 1, X = x] \quad \text{(treated outcome)} \\
e(x) &= P(T = 1 \mid X = x) \quad \text{(propensity score)}
\end{align}
The naive CATE estimator $\hat{\tau}(x) = \hat{\mu}_1(x) - \hat{\mu}_0(x)$ inherits regularization bias from ML estimation.
\end{definition}


% ============================================================================
\section{DML for Heterogeneous Effects}\label{sec:ch09_dml_heterogeneity}
% ============================================================================

The key insight from Chernozhukov et al. (2018) extends naturally to CATE: use Neyman orthogonal scores to eliminate first-stage bias.

\subsection{The R-Learner Framework}

\begin{definition}[R-Learner for CATE]
\label{def:r_learner}
The R-Learner (\cite{nie2021quasi}) minimizes a weighted squared loss:
\begin{equation}
\hat{\tau}(\cdot) = \argmin_{\tau \in \mathcal{T}} \sum_{i=1}^{n} \left( \underbrace{Y_i - \hat{m}(X_i)}_{\text{residualized } Y} - \underbrace{(T_i - \hat{e}(X_i))}_{\text{residualized } T} \cdot \tau(X_i) \right)^2
\label{eq:r_learner}
\end{equation}
where $\hat{m}(x) = \mathbb{E}[Y \mid X = x]$ (marginal outcome).
\end{definition}

This is exactly DML applied point-wise:
\begin{enumerate}
    \item Residualize $Y$ on $X$: $\tilde{Y}_i = Y_i - \hat{m}(X_i)$
    \item Residualize $T$ on $X$: $\tilde{T}_i = T_i - \hat{e}(X_i)$
    \item Regress $\tilde{Y}$ on $\tilde{T} \cdot \tau(X)$
\end{enumerate}

\subsection{Implementation with EconML}

EconML provides production-ready implementations:

\begin{minted}[fontsize=\small]{python}
from econml.dml import LinearDML, CausalForestDML
from sklearn.ensemble import RandomForestRegressor, RandomForestClassifier

# Generate heterogeneous DGP
np.random.seed(42)
n = 2000
X = np.random.randn(n, 5)
T = np.random.binomial(1, 0.5, n)

# True heterogeneous effect: tau(X) = 1 + 2*X_0
true_tau = 1 + 2 * X[:, 0]
Y = true_tau * T + X @ np.array([1, 0.5, -0.3, 0.2, 0.1]) + np.random.randn(n)

# LinearDML with effect modifiers
est_linear = LinearDML(
    model_y=RandomForestRegressor(n_estimators=100, n_jobs=-1),
    model_t=RandomForestClassifier(n_estimators=100, n_jobs=-1),
    discrete_treatment=True,
    cv=5
)
est_linear.fit(Y, T, X=X[:, :2], W=X[:, 2:])  # X=effect modifiers, W=controls

# Get CATE estimates
cate_linear = est_linear.effect(X[:, :2])
\end{minted}

\insight{
\textbf{X vs W in EconML}

EconML distinguishes:
\begin{itemize}
    \item \textbf{X (effect modifiers)}: Variables the CATE depends on---$\tau(X)$
    \item \textbf{W (controls)}: Variables needed for unconfoundedness but not effect heterogeneity
\end{itemize}

Putting everything in X leads to curse of dimensionality. Putting effect modifiers in W throws away heterogeneity information. Good practice: domain knowledge guides the split.
}


% ============================================================================
\section{CausalForestDML}\label{sec:ch09_causal_forest}
% ============================================================================

When effect heterogeneity is non-linear or involves interactions, parametric models fail. CausalForestDML combines random forests with DML orthogonalization.

\subsection{Theory}

\begin{definition}[Causal Forest]
\label{def:causal_forest}
A causal forest (\cite{wager2018estimation}) is a random forest where:
\begin{enumerate}
    \item Each tree splits to maximize heterogeneity in $\tau(X)$
    \item Predictions are \textbf{honest}: training data for splits differs from prediction data
    \item Local estimates use a neighborhood weighted by tree proximity
\end{enumerate}
The splitting criterion maximizes:
\begin{equation}
\text{Heterogeneity gain} \propto (\hat{\tau}_{\text{left}} - \hat{\tau}_{\text{right}})^2
\end{equation}
\end{definition}

CausalForestDML adds DML orthogonalization: first residualize $Y$ and $T$, then fit a causal forest to the residuals.

\subsection{Implementation}

\begin{minted}[fontsize=\small]{python}
from econml.dml import CausalForestDML
from sklearn.ensemble import RandomForestRegressor, RandomForestClassifier

# CausalForestDML for non-parametric heterogeneity
est_cf = CausalForestDML(
    model_y=RandomForestRegressor(n_estimators=100, n_jobs=-1),
    model_t=RandomForestClassifier(n_estimators=100, n_jobs=-1),
    discrete_treatment=True,
    n_estimators=200,        # Number of trees in causal forest
    min_samples_leaf=10,     # Minimum samples per leaf
    max_depth=None,          # Grow trees to full depth
    honest=True,             # Use honest trees (separate split/estimation data)
    inference=True,          # Enable confidence intervals
    cv=5,
    n_jobs=-1,
    random_state=42
)

est_cf.fit(Y, T, X=X, W=None)

# CATE predictions with confidence intervals
cate_pred = est_cf.effect(X)
cate_lower, cate_upper = est_cf.effect_interval(X, alpha=0.05)

print(f"Average CATE: {cate_pred.mean():.3f}")
print(f"CATE std: {cate_pred.std():.3f}")
print(f"True tau std: {true_tau.std():.3f}")
\end{minted}

\subsection{Inference}

CausalForestDML provides valid inference through the \texttt{effect\_inference} method:

\begin{minted}[fontsize=\small]{python}
# Detailed inference for individual units
inference_results = est_cf.effect_inference(X[:5])
summary_df = inference_results.summary_frame(alpha=0.05)
print(summary_df)

# Columns: point_estimate, stderr, zstat, pvalue, ci_lower, ci_upper

# Population summary
pop_summary = est_cf.effect_inference(X).population_summary(alpha=0.05)
print(pop_summary)
\end{minted}


% ============================================================================
\section{Best Linear Predictor of CATE}\label{sec:ch09_blp}
% ============================================================================

While $\hat{\tau}(x)$ gives individual predictions, the \textbf{Best Linear Predictor} (BLP) summarizes heterogeneity in an interpretable way.

\subsection{Theory}

\begin{definition}[Best Linear Predictor]
\label{def:blp}
The BLP projects the true CATE onto a linear function of effect modifiers:
\begin{equation}
\tau_{\text{BLP}}(X) = \alpha + X'\beta^*
\end{equation}
where $\beta^* = \argmin_{\beta} \mathbb{E}[(\tau(X) - X'\beta)^2]$.

The BLP answers: ``Which covariates \textit{linearly} predict effect heterogeneity?''
\end{definition}

\begin{theorem}[BLP Interpretation]
\label{thm:blp_interpretation}
The BLP coefficient $\beta^*_j$ measures the \textbf{marginal association} between covariate $X_j$ and treatment effect $\tau(X)$, holding other covariates fixed (in a linear projection sense).
\end{theorem}

\subsection{Implementation}

EconML's \texttt{const\_marginal\_effect} and model summaries provide BLP-like quantities:

\begin{minted}[fontsize=\small]{python}
from econml.dml import LinearDML

# LinearDML naturally gives BLP interpretation
est_linear = LinearDML(
    model_y=RandomForestRegressor(n_estimators=100, n_jobs=-1),
    model_t=RandomForestClassifier(n_estimators=100, n_jobs=-1),
    discrete_treatment=True,
    fit_cate_intercept=True,
    cv=5
)

# Effect modifiers: X[:, :2], Controls: X[:, 2:]
est_linear.fit(Y, T, X=X[:, :2], W=X[:, 2:])

# BLP coefficients (effect of X on tau)
print("\n=== Best Linear Predictor Coefficients ===")
print(est_linear.summary())

# The const_marginal_effect gives the intercept (ATE)
ate = est_linear.const_marginal_effect()
print(f"\nATE (CATE intercept): {ate[0]:.3f}")
\end{minted}

\subsection{Chernozhukov et al. (2018) BLP Test}

The original DML paper proposes a specific BLP analysis:

\begin{minted}[fontsize=\small]{python}
def compute_blp_statistics(est, X_test, y_test, t_test, w_test=None):
    """
    Compute BLP statistics following Chernozhukov et al. (2018).

    Returns:
        - beta_1: BLP coefficient (measures heterogeneity)
        - standard error
        - p-value for testing H0: beta_1 = 0 (no heterogeneity)
    """
    # Get CATE predictions
    cate_pred = est.effect(X_test)

    # Center the CATE predictions
    cate_centered = cate_pred - cate_pred.mean()

    # BLP regression: tau_hat = alpha + beta * (tau_tilde - mean(tau_tilde))
    # Under the null of no heterogeneity, beta = 0
    from scipy import stats
    import statsmodels.api as sm

    # Use weighted least squares with propensity weights
    X_blp = sm.add_constant(cate_centered)
    model = sm.OLS(cate_pred, X_blp)
    results = model.fit()

    beta_1 = results.params[1]
    se_beta_1 = results.bse[1]
    pval = results.pvalues[1]

    return {
        'beta_1': beta_1,
        'se': se_beta_1,
        'pvalue': pval,
        'significant_heterogeneity': pval < 0.05
    }

blp_stats = compute_blp_statistics(est_cf, X, Y, T)
print(f"BLP beta_1: {blp_stats['beta_1']:.3f}")
print(f"P-value for heterogeneity: {blp_stats['pvalue']:.4f}")
\end{minted}


% ============================================================================
\section{Visualization and Interpretation}\label{sec:ch09_visualization}
% ============================================================================

Effective visualization reveals heterogeneity patterns.

\subsection{CATE Distribution}

\begin{minted}[fontsize=\small]{python}
import matplotlib.pyplot as plt

fig, axes = plt.subplots(1, 3, figsize=(14, 4))

# Distribution of CATE estimates
ax = axes[0]
ax.hist(cate_pred, bins=30, edgecolor='black', alpha=0.7)
ax.axvline(cate_pred.mean(), color='red', linestyle='--',
           label=f'Mean={cate_pred.mean():.2f}')
ax.axvline(0, color='black', linestyle='-', alpha=0.5)
ax.set_xlabel('Estimated CATE')
ax.set_ylabel('Frequency')
ax.set_title('Distribution of Treatment Effects')
ax.legend()

# CATE vs key covariate
ax = axes[1]
sort_idx = np.argsort(X[:, 0])
ax.scatter(X[sort_idx, 0], cate_pred[sort_idx], alpha=0.3, s=10)
ax.plot(X[sort_idx, 0], true_tau[sort_idx], 'r-',
        linewidth=2, label='True CATE')
ax.set_xlabel('X_0 (key effect modifier)')
ax.set_ylabel('CATE')
ax.set_title('CATE vs Effect Modifier')
ax.legend()

# Confidence intervals for sorted CATE
ax = axes[2]
sorted_idx = np.argsort(cate_pred)[:100]  # Bottom 100
x_plot = np.arange(len(sorted_idx))
ax.errorbar(x_plot, cate_pred[sorted_idx],
            yerr=[cate_pred[sorted_idx] - cate_lower[sorted_idx],
                  cate_upper[sorted_idx] - cate_pred[sorted_idx]],
            fmt='o', markersize=3, capsize=2, alpha=0.5)
ax.axhline(0, color='red', linestyle='--')
ax.set_xlabel('Unit (sorted by CATE)')
ax.set_ylabel('CATE with 95% CI')
ax.set_title('Individual CATE Estimates')

plt.tight_layout()
plt.savefig('cate_visualization.pdf', bbox_inches='tight')
\end{minted}

\subsection{Partial Dependence Plots}

\begin{minted}[fontsize=\small]{python}
def plot_cate_partial_dependence(est, X, feature_idx, feature_name, n_grid=50):
    """
    Plot CATE as a function of one feature, averaging over others.

    This is analogous to partial dependence plots for ML models.
    """
    x_grid = np.linspace(X[:, feature_idx].min(),
                         X[:, feature_idx].max(), n_grid)

    cate_grid = []
    cate_lower_grid = []
    cate_upper_grid = []

    for x_val in x_grid:
        # Create modified X with feature_idx fixed at x_val
        X_mod = X.copy()
        X_mod[:, feature_idx] = x_val

        # Average CATE over all other covariates
        cate = est.effect(X_mod).mean()
        lower, upper = est.effect_interval(X_mod, alpha=0.05)

        cate_grid.append(cate)
        cate_lower_grid.append(lower.mean())
        cate_upper_grid.append(upper.mean())

    plt.figure(figsize=(8, 5))
    plt.plot(x_grid, cate_grid, 'b-', linewidth=2, label='CATE')
    plt.fill_between(x_grid, cate_lower_grid, cate_upper_grid,
                     alpha=0.2, label='95% CI')
    plt.axhline(0, color='red', linestyle='--', alpha=0.5)
    plt.xlabel(feature_name)
    plt.ylabel('Conditional Average Treatment Effect')
    plt.title(f'CATE Partial Dependence: {feature_name}')
    plt.legend()
    plt.tight_layout()

    return x_grid, cate_grid

# Example usage
plot_cate_partial_dependence(est_cf, X, 0, 'X_0 (effect modifier)')
\end{minted}


% ============================================================================
\section{Subgroup Analysis and Policy Trees}\label{sec:ch09_subgroups}
% ============================================================================

Heterogeneity analysis often culminates in \textbf{actionable subgroups}---identifying who should receive treatment.

\subsection{CATE-Based Subgroups}

\begin{minted}[fontsize=\small]{python}
def analyze_subgroups(cate_pred, X, feature_names, n_groups=4):
    """
    Analyze treatment effect heterogeneity by CATE quartiles.
    """
    quartiles = np.percentile(cate_pred, [25, 50, 75])

    groups = np.digitize(cate_pred, quartiles)

    results = []
    for g in range(4):
        mask = groups == g
        group_data = {
            'group': g + 1,
            'n': mask.sum(),
            'cate_mean': cate_pred[mask].mean(),
            'cate_std': cate_pred[mask].std(),
        }

        # Add covariate means for interpretation
        for i, name in enumerate(feature_names):
            group_data[f'{name}_mean'] = X[mask, i].mean()

        results.append(group_data)

    import pandas as pd
    return pd.DataFrame(results)

feature_names = [f'X_{i}' for i in range(X.shape[1])]
subgroups = analyze_subgroups(cate_pred, X, feature_names)
print(subgroups.to_string())
\end{minted}

\subsection{Policy Trees with EconML}

EconML provides \texttt{SingleTreeCateInterpreter} for interpretable subgroup identification:

\begin{minted}[fontsize=\small]{python}
from econml.cate_interpreter import SingleTreeCateInterpreter

# Fit interpretable tree to CATE
intrp = SingleTreeCateInterpreter(
    include_model_uncertainty=True,
    max_depth=3,               # Shallow for interpretability
    min_samples_leaf=50,       # Ensure reliable subgroups
    min_impurity_decrease=0.01
)

# X should be the effect modifiers used in estimation
intrp.interpret(est_cf, X)

# Visualize the tree
plt.figure(figsize=(20, 10))
intrp.plot(feature_names=feature_names, fontsize=10)
plt.title('CATE Interpretation Tree')
plt.tight_layout()
plt.savefig('cate_tree.pdf', bbox_inches='tight')
\end{minted}

\subsection{Policy Trees for Treatment Assignment}

Beyond interpretation, we can learn \textbf{optimal treatment policies}:

\begin{minted}[fontsize=\small]{python}
from econml.cate_interpreter import SingleTreePolicyInterpreter

# Policy tree: who should be treated?
# Accounts for treatment costs
policy_intrp = SingleTreePolicyInterpreter(
    include_model_uncertainty=True,
    max_depth=2,
    min_samples_leaf=100,
    min_impurity_decrease=0.001
)

# Treatment cost: only treat if CATE > 0.2
treatment_cost = 0.2 * np.ones(len(X))
policy_intrp.interpret(est_cf, X, sample_treatment_costs=treatment_cost)

# Visualize policy
plt.figure(figsize=(15, 8))
policy_intrp.plot(feature_names=feature_names, fontsize=12)
plt.title('Optimal Treatment Policy Tree (cost = 0.2)')
plt.tight_layout()
\end{minted}

\insight{
\textbf{Policy Trees vs. Prediction Trees}

\textbf{Prediction trees} minimize squared error in predicting $\tau(x)$. They may include irrelevant splits if they marginally improve prediction.

\textbf{Policy trees} maximize \textit{policy value}---the expected benefit from following the tree's treatment rule. Splits that don't change the optimal decision are pruned.

For deployment, policy trees provide cleaner, more actionable rules.
}


% ============================================================================
\section{Application: Insurance Pricing Heterogeneity}\label{sec:ch09_application}
% ============================================================================

Let's apply heterogeneity analysis to the insurance pricing context from Chapter 8.

\begin{minted}[fontsize=\small]{python}
from src.validation import create_insurance_dgp
from econml.dml import CausalForestDML

# Generate insurance DGP with heterogeneous effects
dgp = create_insurance_dgp(
    realism="moderate",
    n_periods=120,
    n_products=20,
    true_tau=-0.8,  # Base effect
    seed=42
)

# Effect modifiers: product characteristics that might affect tau
# - product_size: larger products may be less price-sensitive
# - distribution_channel: different channels have different elasticities
# - customer_age: demographics affect price sensitivity

# Simulate heterogeneity based on product characteristics
n_obs = len(dgp.Y)
product_idx = dgp.panel_info['product_idx']
n_products = dgp.params.n_products

# Create product-level effect modifiers
np.random.seed(42)
product_size = np.random.exponential(1.0, n_products)
channel_type = np.random.choice([0, 1], n_products)  # 0=direct, 1=broker

X_effect = np.column_stack([
    product_size[product_idx],
    channel_type[product_idx]
])

# Controls: macro confounders
W_controls = dgp.macro_controls

# Fit CausalForestDML
est_insurance = CausalForestDML(
    model_y='auto',
    model_t='auto',
    discrete_treatment=False,
    n_estimators=200,
    min_samples_leaf=20,
    cv=3,
    n_jobs=-1,
    random_state=42
)

est_insurance.fit(dgp.Y, dgp.T, X=X_effect, W=W_controls)

# Analyze heterogeneity
cate_insurance = est_insurance.effect(X_effect)
print(f"ATE: {cate_insurance.mean():.3f}")
print(f"CATE range: [{cate_insurance.min():.3f}, {cate_insurance.max():.3f}]")

# By product size quartile
size_quartiles = np.percentile(X_effect[:, 0], [25, 50, 75])
for q, (low, high) in enumerate(zip([0] + list(size_quartiles),
                                     list(size_quartiles) + [np.inf])):
    mask = (X_effect[:, 0] >= low) & (X_effect[:, 0] < high)
    print(f"Size Q{q+1}: CATE = {cate_insurance[mask].mean():.3f}")
\end{minted}


% ============================================================================
\section{Summary}\label{sec:ch09_summary}
% ============================================================================

This chapter extended DML from average effects to heterogeneous effects:

\begin{enumerate}
    \item \textbf{CATE}: $\tau(x) = \mathbb{E}[Y(1) - Y(0) \mid X = x]$ quantifies individual-level effects
    \item \textbf{DML for CATE}: Orthogonal scores eliminate first-stage bias in CATE estimation
    \item \textbf{CausalForestDML}: Non-parametric heterogeneity discovery with valid inference
    \item \textbf{BLP}: Linear summary of heterogeneity for interpretation and testing
    \item \textbf{Policy trees}: Actionable rules for treatment targeting
\end{enumerate}

\keyresult{
\textbf{Key Takeaways}
\begin{itemize}
    \item Heterogeneity reveals who benefits from treatment---essential for targeting
    \item CausalForestDML combines random forests with DML orthogonalization
    \item EconML's X/W split distinguishes effect modifiers from controls
    \item BLP tests whether heterogeneity is statistically significant
    \item Policy trees translate CATE into actionable treatment rules
\end{itemize}
}


% ============================================================================
\section{Exercises}\label{sec:ch09_exercises}
% ============================================================================

\begin{enumerate}
    \item \textbf{CATE Recovery}: Generate data with known heterogeneity $\tau(x) = 1 + 2x_1 - 0.5x_1^2$. Compare LinearDML and CausalForestDML recovery of this non-linear pattern.

    \item \textbf{BLP Testing}: Using the 401(k) data from Chapter 3, test for heterogeneity by income and age. Is there statistically significant effect heterogeneity?

    \item \textbf{Policy Optimization}: For the insurance pricing application, construct a policy tree that maximizes expected profit (CATE minus treatment cost). How does the optimal policy change with cost assumptions?

    \item \textbf{Time-Varying Heterogeneity}: Extend the rolling window analysis from Chapter 6 to estimate time-varying CATE. Does heterogeneity structure change over time?

    \item \textbf{Subgroup Stability}: Using bootstrap resampling, assess the stability of the top CATE quartile. What fraction of units are consistently in the highest-effect group?
\end{enumerate}
